\documentclass[12pt]{article}

%%%%%%%%%%%%%%%%%%%%%%%%%%%%
%%%%%%%%%%%%%%%%%%%%%%%%%%%%
% Load in packages
\usepackage{amsmath}
\usepackage{float}
\usepackage{amssymb}
\usepackage{hyperref}
\usepackage{multicol}
\usepackage{multirow}
\usepackage{graphicx}
\usepackage{enumitem}
\usepackage[top=1in, bottom=1in, left=1in, right=1in]{geometry}

%%%%%%%%%%%%%%%%%%%%%%%%%%%%
%%%%%%%%%%%%%%%%%%%%%%%%%%%%

\begin{document}

\begin{center}
\Large Chapter 17 Practice Problems

\medskip

\normalsize Elements of Microeconomics (discussion section 1)

\medskip

\small Rudy Kretschmar
\end{center}

\section*{Question 1}
Consider a monopolistically competitive market with $N$ firms in it. The following demand, marginal revenue, total cost, and marginal cost equations are for a single firm in the market:
\[Q_D = 50/N - 2P\]
\[MR = 100/N - Q\]
\[TC = 30 + Q^2\]
\[MC = 2Q\]
\begin{enumerate}[label = (\alph*)]
    \item How does the number of firms in the market ($N$) affect the demand curve of each firm?
    \item At what $Q$ will each firm produce? 
    \item What price does each firm charge?
    \item How much profit does each firm make?
    \item In the long run (where entry and exit are possible), how many firms will there be in this market?
\end{enumerate}

\section*{Question 2}
Consider the market for grocery stores. Giant is one of the many firms in this masrket, which is in long-run equilibrium.
\begin{enumerate}[label = (\alph*)]
    \item Draw a diagram showing Giant's demand curve, marginal revenue, marginal cost, and average total cost. Label the profit maximizing  $Q$ and $P$ for Giant.
    \item What is the efficient/socially optimal level of output and price in this market? Label this on your diagram.
    \item What is Giant's profit? Why?
\end{enumerate}

\newpage
\section*{Question 3}
How would you classify each of the following markets, perfect competition, monopolistic competition, or monopoly? Explain why.
\begin{enumerate}[label = (\alph*)]
    \item Tap water
    \item Lipstick
    \item Local electricity
    \item Peanut butter
    \item Train tickets
\end{enumerate}

\section*{Question 4}
In Baltimore there are two main internet providers, Xfinity and Verizon. Suppose the marginal cost of supplying internet is constant at \$20 per customer, and the demand for internet service is described by the following demand schedule:
\begin{table}[H]
    \centering
    \begin{tabular}{c|p}
        Price & Quantity \\
        \hline
        \$200 & 5,000 customers \\
        \$180 & 6,000 \\
        \$160 & 7,000 \\
        \$140 & 8,000 \\
        \$120 & 9,000 \\
        \$100 & 10,000 \\
        \$80 & 11,000 \\
        \$60 & 12,000 \\
    \end{tabular}
    \label{tab:placeholder}
\end{table}
\begin{enumerate}[label = (\alph*)]
    \item If there were many suppliers of internet service in Baltimore, what kind of market would this be? What would be the price and quantity? 
    \item If there were only one supplier of internet service, what kind of market would this be? What would the price and quantity? 
    \item If Xfinity and Verizon form a cartel, what would be the price and quantity?
    \begin{enumerate}[label = (\roman*)]
        \item If Xfinity and Verizon split the market evenly in this case, what would be Verizon's production and profit?
        \item What would happen to Verizon's profit if it increased it's production by 1,000 while Xfinity stuck with the cartel agreement?
    \end{enumerate}
    \item Why are cartel agreements often unsuccessful?
\end{enumerate}

\newpage

\section*{Question 5}
Sandy's Brew is a small coffee company that is considering entering the local coffee market which is dominated by Common Ground Coffee. Each company's profit depends on whether Sandy's Brew enters the market and whether Common Ground Coffee sets a hgih price or a low price:

\begin{center}
\begin{tabular}{c|c|c|c|}
   & & \multicolumn{2}{c|}{\textbf{Common Ground Coffee}} \\
   & & High Price & Low Price \\
    \hline
    \multirow{2}{*}{\textbf{Sandy's Brew}} & Enter & $(\$500,\$750)$ & $(-\$250,\$250)$ \\
    \cline{2-4}
    & Don't Enter & $(\$0,\$1600)$ & $(\$0,\$500)$ \\
    \hline
\end{tabular}
\end{center}

Where the row player earns the first value in the parentheses in each situation, and the column player earns the second value in the parentheses in each situation.
\begin{enumerate}[label = (\alph*)]
    \item Does either player have a dominant strategy?
    \item What is a nash equilibrium?
    \item Is there a nash equilibrium in this game? If so what is it?
    \begin{enumerate}[label = (\roman*)]
        \item How does your answer to (a) help you determine the nash equilibrium
    \end{enumerate}
    \item Common Ground Coffee threatens to set a low price if Sandy's Brew enters the market. Is this a credible threat? Why or why not?
\end{enumerate}

\
\section*{Question 6}
Consider the market for cigarettes, which has a demand and supply curve given by:
\[Q_D = 75 - 3P\]
\[Q_S = 5 + P\]
\begin{enumerate}[label = (\alph*)]
    \item What is the equilibrium price and quantity in this market?
    \item Cigarettes have a negative externality on bystanders due to second hand smoke, causing the socially optimal quantity of cigarettes to be less than the equilibrium quantity in this market. Through what intervention can the government bring the market equilibrium closer to the socially optimal level?
    \item If the socially optimal quantity of cigarettes is 20, a tax of what amount per pack of cigarettes must be levied by the government to bring the market equilibrium to equal the the socially optimal quantity.
\end{enumerate}

\end{document}
